\documentclass[tikz,border=8pt]{standalone}
\usepackage{tikz}
\usepackage{amsmath}
\usepackage{pgfplots}
\pgfplotsset{compat=1.18}
\usepackage{ctex}
\usetikzlibrary{calc, positioning, arrows.meta, shapes, fit, backgrounds}

\begin{document}
\begin{tikzpicture}

% ============================================================
% 左子图：Metropolis-Hastings 链轨迹图
% ============================================================
\begin{axis}[
  name=axMH,
  width=8cm, height=5.5cm,
  xmin=0, xmax=1000,
  ymin=-1.0, ymax=4.5,
  xlabel={迭代次数},
  ylabel={状态值 $\theta$},
  axis lines=left,
  xtick={0,200,400,600,800,1000},
  xticklabel style={font=\scriptsize},
  yticklabel style={font=\scriptsize},
  ytick={-1,0,1,2,3,4},
  grid=major, grid style={gray!20},
  clip=true
]

% Burn-in 阶段背景（灰色阴影）
\addplot[fill=gray!20, draw=none] coordinates
  {(0,-1.0) (200,-1.0) (200,4.5) (0,4.5)} \closedcycle;

% Burn-in 阶段轨迹（灰色）
\addplot[gray!60, thick, line join=round] coordinates {
  (0,0.1) (20,0.5) (40,1.2) (60,0.8) (80,1.5) (100,2.0)
  (120,1.6) (140,2.2) (160,1.9) (180,2.3) (200,2.1)
};

% 平稳阶段轨迹（彩色）
\addplot[blue!70!black, thick, line join=round] coordinates {
  (200,2.1) (220,1.8) (240,2.4) (260,2.0) (280,2.5) (300,1.9)
  (320,2.2) (340,2.7) (360,2.0) (380,2.3) (400,1.7)
  (420,2.4) (440,2.1) (460,2.6) (480,2.0) (500,2.3)
  (520,1.8) (540,2.5) (560,2.1) (580,2.4) (600,2.0)
  (620,2.3) (640,1.9) (660,2.5) (680,2.2) (700,1.8)
  (720,2.4) (740,2.1) (760,2.6) (780,2.0) (800,2.3)
  (820,1.9) (840,2.4) (860,2.2) (880,2.5) (900,2.0)
  (920,2.3) (940,1.8) (960,2.4) (980,2.1) (1000,2.2)
};

% 真值水平线
\addplot[red!70!black, dashed, thick, domain=0:1000, samples=2] {2.0};
\node[font=\scriptsize, red!80!black] at (axis cs:850, 2.18) {真值 $\theta^*=2.0$};

% burn-in 标注
\node[font=\scriptsize, gray!70!black] at (axis cs:100, 4.2) {Burn-in};

\end{axis}
\node[font=\normalsize\bfseries] at ($(axMH.north)+(0,0.4cm)$) {Metropolis-Hastings 链};

% ============================================================
% 右子图：Gibbs 采样 二维联合轨迹图
% ============================================================
\begin{axis}[
  name=axGibbs,
  at={(axMH.east)}, anchor=west, xshift=1.2cm,
  width=5.5cm, height=5.5cm,
  xmin=-0.3, xmax=3.8,
  ymin=-0.3, ymax=3.8,
  xlabel={$\theta_1$},
  ylabel={$\theta_2$},
  axis lines=left,
  xtick={0,1,2,3},
  ytick={0,1,2,3},
  xticklabel style={font=\scriptsize},
  yticklabel style={font=\scriptsize},
  grid=major, grid style={gray!20},
  clip=true
]

% Gibbs 轨迹（交替固定一维的折线）
\addplot[orange!70!black, thick, line join=round, mark=none] coordinates {
  (0.2, 0.3)
  (2.1, 0.3)
  (2.1, 1.8)
  (1.0, 1.8)
  (1.0, 2.4)
  (2.3, 2.4)
  (2.3, 1.5)
  (1.5, 1.5)
  (1.5, 2.0)
  (2.0, 2.0)
  (2.0, 1.8)
  (1.8, 1.8)
  (1.8, 2.1)
  (2.1, 2.1)
  (2.1, 1.9)
  (1.9, 1.9)
  (1.9, 2.0)
  (2.0, 2.0)
};

% 起点标记
\addplot[green!50!black, only marks, mark=*, mark size=3pt] coordinates {(0.2, 0.3)};
\node[font=\scriptsize, green!50!black] at (axis cs:0.5, 0.5) {起点};

% 收敛区域（散点群）
\addplot[red!60!black, only marks, mark=o, mark size=1.5pt, opacity=0.6] coordinates {
  (1.9,2.0) (2.0,1.9) (2.1,2.1) (1.8,2.0) (2.0,2.1) (2.2,1.8) (1.9,1.9)
  (2.1,2.0) (2.0,2.2) (1.7,2.0) (2.0,1.8) (1.9,2.1)
};

% 真值标记
\addplot[black, only marks, mark=+, mark size=5pt, thick] coordinates {(2.0, 2.0)};
\node[font=\scriptsize] at (axis cs:2.35, 2.35) {真值};

\end{axis}
\node[font=\normalsize\bfseries] at ($(axGibbs.north)+(0,0.4cm)$) {Gibbs 采样轨迹};

% 整体标题
\node[font=\large\bfseries] at ($(axMH.north)!0.5!(axGibbs.north)+(0,1.3cm)$)
  {MCMC 算法：Metropolis-Hastings 与 Gibbs 采样};

% 底部说明框
\node[font=\small, align=left,
      fill=white, draw=gray!50, thin, rounded corners=2pt, inner sep=5pt]
  at ($(axMH.south)!0.5!(axGibbs.south)-(0,1.0cm)$)
  {MCMC：构造马氏链，使平稳分布为目标分布 $\pi(\theta)$，经充足迭代后样本近似来自 $\pi(\theta)$};

\end{tikzpicture}
\end{document}
